\DocumentMetadata{
  pdfversion=2.0,
  pdfstandard=ua-2,
  lang=en-US,
  tagging=on,
  tagging-setup={math/setup=mathml-SE,tabsorder=structure}
}
\documentclass[11pt,letterpaper]{article}
\usepackage[margin=0.68in,includefoot]{geometry}
\usepackage{newcomputermodern}
\usepackage{amsmath,amscd,mathtools}
\providecommand{\square}{\mdlgwhtsquare}
\usepackage[hidelinks]{hyperref}
\hypersetup{
  pdftitle={Differential Geometry PhD Exam, May 2007},
  pdfauthor={University of Florida Department of Mathematics},
  pdfsubject={Graduate mathematics examination},
  pdfdisplaydoctitle=true,
  bookmarksopen=true
}
\setcounter{secnumdepth}{0}
\setlength{\parindent}{0pt}
\setlength{\parskip}{0.28em}
\setlength{\emergencystretch}{3em}
\setlength{\leftmargini}{2.15em}
\setlength{\leftmarginii}{2.65em}
\setlength{\leftmarginiii}{3em}
\renewcommand{\labelenumi}{\textbf{\theenumi.}}
\pagestyle{empty}
\raggedbottom
\allowdisplaybreaks
\newenvironment{examproblems}[1]{%
  \begin{enumerate}%
  \setcounter{enumi}{#1}%
  \setlength{\topsep}{0.35em}%
  \setlength{\partopsep}{0pt}%
  \setlength{\itemsep}{0.48em}%
  \setlength{\parsep}{0.18em}%
}{\end{enumerate}}
\newenvironment{examsubparts}{%
  \begin{enumerate}%
  \setlength{\topsep}{0.18em}%
  \setlength{\partopsep}{0pt}%
  \setlength{\itemsep}{0.16em}%
  \setlength{\parsep}{0.1em}%
}{\end{enumerate}}
\newenvironment{examinstructions}{%
  \par\begingroup\itshape
}{%
  \par\endgroup\vspace{0.18em}
}
\makeatletter
\renewcommand\section{\@startsection{section}{1}{0pt}%
  {-1.25ex plus -.25ex minus -.1ex}%
  {0.55ex plus .1ex}%
  {\normalfont\large\bfseries}}
\renewcommand{\@maketitle}{%
  \newpage\null\vskip -1.2em%
  {\footnotesize\bfseries
    \href{https://www.ufl.edu/}{University of Florida}, \href{https://math.ufl.edu/}{Department of Mathematics}\par}%
  \vskip 0.18em%
  \tagstructbegin{tag=Title}%
    {\LARGE\bfseries\href{https://gma.math.ufl.edu/exams/differential-geometry-phd/}{Differential Geometry PhD Exam}, May 2007\par}%
  \tagstructend%
  \vskip 0.35em%
  \tagmcbegin{artifact}%
    \hrule height 1pt%
  \tagmcend%
  \vskip 0.55em%
}
\makeatother
\title{Differential Geometry PhD Exam, May 2007}
\author{}
\date{}
\begin{document}
\maketitle
\thispagestyle{empty}
\begin{examinstructions}
Attempt SIX problems. Write solutions in a neat and logical fashion, giving complete reasons for all steps and stating carefully any substantial theorems used.
\end{examinstructions}
\begin{examproblems}{0}
\item
Decide which of the following are submanifolds of \(\mathbb{R}^3\):
\begin{examsubparts}
\item[\textnormal{(a)}]
\(A=\{(x,y,z):z^2=x^2+y^2-1\}\);
\item[\textnormal{(b)}]
\(B=\{(x,y,z):z^2=x^2+y^2+1\}\);
\item[\textnormal{(c)}]
\(C=\{(x,y,z):z^2=x^2+y^2\}\).
\end{examsubparts}
\item
Defining the terms involved, prove the Cartan identity 
\[
L_\zeta=(\zeta\mathbin{\lrcorner})\circ d+d\circ(\zeta\mathbin{\lrcorner})
\]
 for the Lie derivative of differential forms along the vector field~\(\zeta\).
\item
Let \(\theta\in\Omega^1(M)\) be nonsingular (that is, nowhere zero) with \(\theta\wedge d\theta=0\). Briefly explaining why it is possible, choose \(\alpha\in\Omega^1(M)\) such that \(d\theta=\theta\wedge\alpha\) and choose \(\beta\in\Omega^1(M)\) such that \(d\alpha=\theta\wedge\beta\).
\begin{examsubparts}
\item[\textnormal{(a)}]
Prove that the three-form \(\alpha\wedge d\alpha\) on~\(M\) is closed.
\item[\textnormal{(b)}]
Prove that if also \(d\theta=\theta\wedge\alpha'\) then \(\alpha'\wedge d\alpha'-\alpha\wedge d\alpha\) is exact.
\end{examsubparts}
\item
Let \(k\in\mathbb{R}\) and define \(\omega_m\in\Omega^m(\mathbb{R}^{m+1}-\{0\})\) by 
\[
(x_0^2+\cdots+x_m^2)^k\omega_m=\sum_{j=0}^m(-1)^j x_j\,dx_0\wedge\cdots\wedge\widehat{dx_j}\wedge\cdots\wedge dx_m,
\]
 where the circumflex over a term indicates that it is omitted.
\begin{examsubparts}
\item[\textnormal{(a)}]
Determine the value(s) of~\(k\) (if any) for which the~\(m\)-form \(\omega_m\) is closed.
\item[\textnormal{(b)}]
Determine the value(s) of~\(k\) (if any) for which the~\(2\)-form \(\omega_2\) is exact.
\end{examsubparts}
\item
Define the Poisson bracket between smooth functions on a symplectic manifold, both in invariant form and in terms of (local) Darboux coordinates. Let \(\omega=\sum_{j=1}^3dp_j\wedge dq_j\) be the standard symplectic form on \(\mathbb{R}^6\) and define 
\[
L_1=q_2p_3-q_3p_2,\qquad L_2=q_3p_1-q_1p_3,\qquad L_3=q_1p_2-q_2p_1.
\]
 Explaining the terms, prove that if \(L_1\) and \(L_2\) are constants of the motion for a given Hamiltonian function on \((\mathbb{R}^6,\omega)\) then so is \(L_3\).
\item
Let~\(G\) be both a group and a smooth manifold; suppose that the group operation \(p:G\times G\to G\) is smooth. By considering the function 
\[
f:G\times G\to G\times G:(x,y)\mapsto(x,p(x,y))
\]
 or otherwise, prove that the inversion map \(G\to G:g\mapsto g^{-1}\) is smooth. [It may be assumed that the tangent map of~\(p\) at \((a,b)\) is given by 
\[
\xi\in T_aG,\ \eta\in T_bG\Rightarrow p_*(\xi,\eta)=\lambda_*^a\eta+\rho_*^b\xi,
\]
 where \(\lambda^a:G\to G\times G:y\mapsto(a,y)\) and \(\rho^b:G\to G\times G:x\mapsto(x,b)\).]
\item
Let \(\sigma:G\times M\to M:(g,x)\mapsto\sigma_g(x)\) be a smooth (left) action of the Lie group~\(G\) on the smooth manifold~\(M\). Explain how the induced map \(\dot{\sigma}:\mathfrak{g}\to\mathrm{Vec}\,M\) is defined. For \(\xi\in\mathfrak{g}\) and \(x\in M\) define 
\[
\gamma:\mathbb{R}\to M:t\mapsto\sigma_{\exp(t\xi)}(x).
\]
 Show that the tangent vector to this curve at time~\(t\) is given by 
\[
\dot{\gamma}(t)=(\sigma_{\exp(t\xi)})_*\dot{\gamma}(0).
\]
 Hence show that if the action~\(\sigma\) is free and \(\dot{\sigma}(\xi)_x=0\) for some \(x\in M\), then \(\xi=0\).
\end{examproblems}
\end{document}
